\chapter{Comparación de los cinco casos}
\label{chap:comparativa}

\section{Síntesis de resultados}
La tabla~\ref{tab:comparativa} compara el mecanismo
documentado, las consecuencias y la prioridad defensiva.
Distingue los hechos observados de las limitaciones
de la información pública.

\begin{table}[htbp]
    \centering
    \small
    \renewcommand{\arraystretch}{1.15}
    \begin{tabularx}{\textwidth}{
        @{}p{2.1cm}XXX@{}
        }
        \toprule
        Caso                                  & Mecanismo & Impacto observado & Prioridad defensiva \\
        \midrule
        PHP2Shell                             &
        Fallos de memoria y encadenamiento con escalada
        de privilegios.                       &
        Compromiso de servidores, robo de credenciales
        y acceso interno.                     &
        Actualizar PHP y productos que lo integran;
        investigar accesos previos.                                                                 \\

        Ivanti                                &
        Evasión de autenticación, ejecución de código
        y escalada a \textit{root}.           &
        Extracción de credenciales y metadatos
        del directorio corporativo.           &
        Corregir toda la cadena, realizar análisis
        forense y revocar credenciales.                                                             \\

        Moldavia                              &
        Deficiencias de autenticación y autorización
        en servicios públicos.                &
        Extracción de datos personales e interrupciones
        preventivas.                          &
        Validar permisos por operación y registro;
        detectar consultas anómalas.                                                                \\

        Francia                               &
        Manipulación de controladores de refrigeración;
        vector concreto no acreditado.        &
        Alteración de la cadena de frío
        y pérdida de alimentos.               &
        Retirar la exposición directa y asegurar
        el mantenimiento remoto.                                                                    \\

        Corea del Sur                         &
        Inyección SQL y evasión de autenticación,
        según ThreatLabz, con apoyo de ARTEX. &
        Extracción de información personal
        y financiera.                         &
        Corregir aplicaciones periféricas y
        controlar accesos y extracciones.                                                           \\
        \bottomrule
    \end{tabularx}
    \caption{Comparación de los incidentes analizados.
        Síntesis basada en las fuentes de los capítulos anteriores.}
    \label{tab:comparativa}
\end{table}

\section{Patrones y diferencias}
El primer patrón es la importancia de las interfaces
accesibles desde el exterior. Las VPN, aplicaciones
periféricas y consolas de mantenimiento conectan
necesidades legítimas de acceso con activos sensibles.
Su exposición exige controles acordes con las
consecuencias de un compromiso, especialmente cuando
permiten alcanzar redes internas o modificar procesos
físicos~\cite{volexity2026ivanti,certfr2026frigorificos,
    zscaler2026corea}.

El segundo patrón es que contener el punto de entrada
no revierte todo el daño. En PHP2Shell e Ivanti,
la recuperación debe considerar credenciales y sistemas
alcanzados; en Moldavia y Corea del Sur, la información
extraída puede facilitar fraudes posteriores.
Francia añade una diferencia decisiva: restaurar
el controlador no recupera los alimentos
perdidos~\cite{dream2026php,volexity2026ivanti,
    age2026datos,zscaler2026corea,onyphe2026francia}.

Finalmente, los casos requieren respuestas distintas.
Los fallos conocidos demandan actualización y comprobación
de versiones; las deficiencias de autorización exigen
revisar qué puede consultar cada usuario.
El caso surcoreano incorpora IA a la ejecución ofensiva,
pero las debilidades descritas siguen requiriendo
controles de seguridad de aplicaciones
convencionales~\cite{zscaler2026corea}.

\section{Límites de la comparación}
No se establece una clasificación por CVSS:
PHP2Shell utiliza la versión 4.0 e Ivanti la 3.1,
mientras que otros casos carecen de puntuaciones
confirmadas para el fallo explotado.
Tampoco son equivalentes servidores comprometidos,
registros extraídos, personas afectadas y establecimientos.
La comparación es cualitativa y está limitada por
investigaciones abiertas y niveles desiguales
de detalle técnico.